\documentclass[11pt]{article}

\usepackage[margin=0.8in]{geometry}
\usepackage{amsmath,amssymb}
\usepackage{enumitem}
\usepackage{xcolor}
\usepackage{tcolorbox}
\usepackage{fancyhdr}
\usepackage{lastpage}
\usepackage{hyperref}
\usepackage{microtype}

\definecolor{uwpurple}{HTML}{4B2E83}
\definecolor{uwgold}{HTML}{B7A57A}
\definecolor{softpurple}{HTML}{F4F1FA}
\definecolor{softgold}{HTML}{FBF7EA}
\definecolor{deepgold}{HTML}{8A6F2A}

\tcbuselibrary{breakable,skins}

\pagestyle{fancy}
\fancyhf{}
\lhead{\textbf{MATH 394: Probability I}}
\rhead{\textbf{Final Solutions}}
\cfoot{\thepage\ of \pageref{LastPage}}
\setlength{\headheight}{14pt}

\setlength{\parindent}{0pt}
\setlength{\parskip}{0.45em}

\newcommand{\course}{MATH 394: Probability I}
\newcommand{\instructor}{Arman Jahangiri}
\newcommand{\department}{Department of Mathematics}
\newcommand{\university}{University of Washington}

\newtcolorbox{problemheader}{
    colback=softpurple,
    colframe=uwpurple,
    boxrule=0.8pt,
    arc=2mm,
    left=2mm,
    right=2mm,
    top=1mm,
    bottom=1mm
}

\newtcolorbox{solutionbox}[1][] {
    breakable,
    enhanced,
    colback=softgold,
    colframe=deepgold,
    coltitle=white,
    colbacktitle=uwpurple,
    fonttitle=\bfseries,
    title={Solution #1},
    boxrule=0.85pt,
    arc=2mm,
    left=2.5mm,
    right=2.5mm,
    top=1.5mm,
    bottom=1.5mm,
    attach boxed title to top left={xshift=2.5mm,yshift=-2mm},
    boxed title style={
        colback=uwpurple,
        colframe=uwpurple,
        arc=1.5mm,
        boxrule=0pt,
        left=1.6mm,
        right=1.6mm,
        top=0.6mm,
        bottom=0.6mm
    }
}

\newtcolorbox{finalanswerbox}[1][] {
    colback=white,
    colframe=uwpurple,
    boxrule=0.75pt,
    arc=2mm,
    left=2mm,
    right=2mm,
    top=1mm,
    bottom=1mm,
    title={\textbf{Final Answer #1}},
    coltitle=uwpurple,
    colbacktitle=softpurple,
    fonttitle=\bfseries
}

\begin{document}

% ============================================================
% TITLE PAGE
% ============================================================

\begin{titlepage}
\centering
\vspace*{1.2cm}

{\Large \textbf{\university}}\\
{\large \department}

\vspace{1.4cm}

{\Huge \color{uwpurple}\textbf{Final Examination Solutions}}\\[0.25cm]
{\Large \textbf{\course}}\\[0.25cm]
{\large Summer 2026 \quad | \quad Instructor: \instructor}

\vspace{1cm}

\begin{tcolorbox}[
    width=0.78\textwidth,
    colback=softpurple,
    colframe=uwpurple,
    boxrule=0.9pt,
    arc=3mm
]
\centering
\textbf{Exam date:} August 21, 2026

\textbf{Answering time:} 60 minutes

\textbf{Total:} 60 points
\end{tcolorbox}

\end{titlepage}


% ============================================================
% PROBLEM 1
% ============================================================

\begin{problemheader}
\textbf{Problem 1.} \hfill \textbf{16 points}
\end{problemheader}

Let the joint probability density function of random variables \(X\) and \(Y\)
be given by
\[
f(x,y)
=
\begin{cases}
x^2e^{-x(y+1)}, & x\geq0,\quad y\geq0,\\
0, & \text{elsewhere}.
\end{cases}
\]

The marginal probability density functions are
\[
f_X(x)
=
\begin{cases}
xe^{-x}, & x\geq0,\\
0, & \text{otherwise},
\end{cases}
\]
and
\[
f_Y(y)
=
\begin{cases}
\dfrac{2}{(1+y)^3}, & y\geq0,\\
0, & \text{otherwise}.
\end{cases}
\]

\begin{enumerate}[label=(\alph*),leftmargin=*,itemsep=0.8em]

    \item Calculate
    \[
    \operatorname{Cov}(X,Y).
    \]
    \hfill \textbf{(4 points)}

    \item Calculate
    \[
    \operatorname{Cov}(3X+1,\,2X-4).
    \]
    \hfill \textbf{(3 points)}

    \item Are \(X\) and \(Y\) independent? Justify your answer.
    \hfill \textbf{(2 points)}

    \item Calculate the moment-generating function of \(X\), namely
    \[
    M_X(t)=E(e^{tX}).
    \]
    \hfill \textbf{(3 points)}

    \item Using \(M_X(t)\), calculate \(E(X)\).
    \hfill \textbf{(2 points)}
 

\end{enumerate}

\begin{solutionbox}

\textbf{(a)} Recall that
\[
\operatorname{Cov}(X,Y)
=
E(XY)-E(X)E(Y).
\]

First,
\begin{align*}
E(X)
&=
\int_0^\infty x f_X(x)\,dx\\
&=
\int_0^\infty x^2e^{-x}\,dx\\
&=2.
\end{align*}

Next,
\begin{align*}
E(Y)
&=
\int_0^\infty y\frac{2}{(1+y)^3}\,dy\\
&=1.
\end{align*}

To calculate \(E(XY)\), use the joint density:
\begin{align*}
E(XY)
&=
\int_0^\infty\int_0^\infty
xy\,x^2e^{-x(y+1)}\,dy\,dx\\
&=
\int_0^\infty
x^3e^{-x}
\left(
\int_0^\infty y e^{-xy}\,dy
\right)dx.
\end{align*}

For \(x>0\),
\[
\int_0^\infty y e^{-xy}\,dy
=
\frac{1}{x^2}.
\]

Therefore,
\begin{align*}
E(XY)
&=
\int_0^\infty xe^{-x}\,dx\\
&=1.
\end{align*}

Hence
\[
\operatorname{Cov}(X,Y)
=
1-(2)(1)
=
-1.
\]


\textbf{(b)} Using
\[
\operatorname{Cov}(aX+b,cX+d)
=
ac\,\operatorname{Var}(X),
\]
we obtain
\[
\operatorname{Cov}(3X+1,2X-4)
=
6\operatorname{Var}(X).
\]

Now
\[
E(X)=2,
\]
and
\begin{align*}
E(X^2)
&=
\int_0^\infty x^2 f_X(x)\,dx\\
&=
\int_0^\infty x^3e^{-x}\,dx\\
&=6.
\end{align*}

Thus
\[
\operatorname{Var}(X)
=
E(X^2)-[E(X)]^2
=
6-4
=
2.
\]

Therefore,
\[
\operatorname{Cov}(3X+1,2X-4)
=
6(2)
=
12.
\]


\textbf{(c)} If \(X\) and \(Y\) were independent, then we would have
\[
f_{X,Y}(x,y)=f_X(x)f_Y(y)
\]
on their support.

However,
\[
f_X(x)f_Y(y)
=
xe^{-x}\frac{2}{(1+y)^3},
\]
whereas
\[
f_{X,Y}(x,y)
=
x^2e^{-x(y+1)}.
\]

These are not equal in general. Therefore, \(X\) and \(Y\) are not
independent.

Alternatively, part (a) already gives
\[
\operatorname{Cov}(X,Y)=-1\neq0.
\]
Since independence would imply zero covariance, \(X\) and \(Y\) cannot be
independent.


\textbf{(d)} By definition,
\begin{align*}
M_X(t)
&=
E(e^{tX})\\
&=
\int_0^\infty e^{tx}xe^{-x}\,dx\\
&=
\int_0^\infty xe^{-(1-t)x}\,dx.
\end{align*}

For \(t<1\),
\[
\int_0^\infty xe^{-(1-t)x}\,dx
=
\frac{1}{(1-t)^2}.
\]

Therefore,
\[
M_X(t)
=
\frac{1}{(1-t)^2},
\qquad t<1.
\]


\textbf{(e)} We know that
\[
E(X)=M_X'(0).
\]

Since
\[
M_X(t)=(1-t)^{-2},
\]
we have
\[
M_X'(t)=2(1-t)^{-3}.
\]

Therefore,
\[
E(X)
=
M_X'(0)
=
2.
\]

 

\begin{finalanswerbox}

\[
\text{(a)}\quad \operatorname{Cov}(X,Y)=-1.
\]

\[
\text{(b)}\quad
\operatorname{Cov}(3X+1,2X-4)=12.
\]

\[
\text{(c)}\quad X\text{ and }Y\text{ are not independent.}
\]

\[
\text{(d)}\quad
M_X(t)=\frac{1}{(1-t)^2},
\qquad t<1.
\]

\[
\text{(e)}\quad E(X)=2.
\]
 

\end{finalanswerbox}

\end{solutionbox}


\newpage

% ============================================================
% PROBLEM 2
% ============================================================

\begin{problemheader}
\textbf{Problem 2.} \hfill \textbf{8 points}
\end{problemheader}

Let \(X\) and \(Y\) be two positive independent continuous random variables
with probability density functions \(f_1(x)\) and \(f_2(y)\), respectively.

Find the probability density function of
\[
U=\frac{X}{Y}.
\]

\textit{Hint:} Let \(V=X\). Find the joint probability density function of
\(U\) and \(V\) (Jacobian method), and then calculate the marginal probability density function
of \(U\).

\begin{solutionbox}

Define
\[
U=\frac{X}{Y},
\qquad
V=X.
\]

We first find the inverse transformation.

Since
\[
v=x
\]
and
\[
u=\frac{x}{y},
\]
we obtain
\[
x=v,
\qquad
y=\frac{v}{u}.
\]

Because \(X>0\) and \(Y>0\), the support of \((U,V)\) is
\[
u>0,\qquad v>0.
\]

The Jacobian of the inverse transformation is
\[
\frac{\partial(x,y)}{\partial(u,v)}
=
\begin{vmatrix}
\dfrac{\partial x}{\partial u}
&
\dfrac{\partial x}{\partial v}
\\[2mm]
\dfrac{\partial y}{\partial u}
&
\dfrac{\partial y}{\partial v}
\end{vmatrix}.
\]

Since
\[
x=v,
\qquad
y=\frac vu,
\]
we have
\[
\frac{\partial(x,y)}{\partial(u,v)}
=
\begin{vmatrix}
0 & 1\\[2mm]
-\dfrac{v}{u^2} & \dfrac1u
\end{vmatrix}
=
\frac{v}{u^2}.
\]

Thus
\[
\left|
\frac{\partial(x,y)}{\partial(u,v)}
\right|
=
\frac{v}{u^2}.
\]

Since \(X\) and \(Y\) are independent,
\[
f_{X,Y}(x,y)=f_1(x)f_2(y).
\]

Therefore,
\[
f_{U,V}(u,v)
=
f_1(v)
f_2\left(\frac vu\right)
\frac{v}{u^2},
\qquad u,v>0.
\]

Finally, integrate out \(V\):
\[
f_U(u)
=
\int_0^\infty
f_1(v)
f_2\left(\frac vu\right)
\frac{v}{u^2}\,dv,
\qquad u>0.
\]

For \(u\leq0\),
\[
f_U(u)=0.
\]

\begin{finalanswerbox}

\[
f_U(u)
=
\begin{cases}
\displaystyle
\frac{1}{u^2}
\int_0^\infty
v f_1(v)
f_2\left(\frac vu\right)\,dv,
& u>0,\\[4mm]
0,
& u\leq0.
\end{cases}
\]

\end{finalanswerbox}

\end{solutionbox}


\newpage

% ============================================================
% PROBLEM 3
% ============================================================

\begin{problemheader}
\textbf{Problem 3.} \hfill \textbf{7 points}
\end{problemheader}

Let \(X\) and \(Y\) be i.i.d.\ positive random variables.

For each part below, fill in the appropriate equality or inequality symbol.

Write \(=\) if the two sides are always equal, \(\leq\) if the left-hand side
is less than or equal to the right-hand side but not necessarily equal, and
similarly for \(\geq\). If no relation holds in general, write \(?\).

Justify each answer.

\begin{enumerate}[label=(\alph*),leftmargin=*,itemsep=1.2em]

    \item
    \[
    E(\log X)
    \quad \underline{\hspace{1.2cm}} \quad
    \log(E(X)).
    \]
    \hfill \textbf{(3 points)}

    \item
    \[
    P(X\leq Y)
    \quad \underline{\hspace{1.2cm}} \quad
    P(X\geq Y).
    \]
    \hfill \textbf{(2 points)}

    \item
    \[
    E(XY)
    \quad \underline{\hspace{1.2cm}} \quad
    \sqrt{E(X^2)E(Y^2)}.
    \]
    \hfill \textbf{(2 points)}

\end{enumerate}

\begin{solutionbox}

\textbf{(a)} The function
\[
g(x)=\log x
\]
is concave on \((0,\infty)\).

By Jensen's inequality for a concave function,
\[
E[g(X)]\leq g(E[X]).
\]

Therefore,
\[
E(\log X)\leq\log(E(X)).
\]


\textbf{(b)} Since \(X\) and \(Y\) are i.i.d., the joint distribution of
\((X,Y)\) is symmetric under interchange of \(X\) and \(Y\).

Therefore,
\[
P(X<Y)=P(Y<X).
\]

Also,
\[
P(X\leq Y)
=
P(X<Y)+P(X=Y)
\]
and
\[
P(X\geq Y)
=
P(X>Y)+P(X=Y).
\]

Since
\[
P(X<Y)=P(X>Y),
\]
it follows that
\[
P(X\leq Y)=P(X\geq Y).
\]


\textbf{(c)} By the Cauchy--Schwarz inequality,
\[
|E(XY)|
\leq
\sqrt{E(X^2)E(Y^2)}.
\]

Since \(X\) and \(Y\) are positive,
\[
E(XY)\geq0.
\]

Hence
\[
E(XY)
\leq
\sqrt{E(X^2)E(Y^2)}.
\]

In fact, because \(X\) and \(Y\) are independent,
\[
E(XY)=E(X)E(Y),
\]
but Cauchy--Schwarz gives the requested comparison directly.

\begin{finalanswerbox}

\[
\text{(a)}\quad
E(\log X)
\leq
\log(E(X)).
\]

\[
\text{(b)}\quad
P(X\leq Y)
=
P(X\geq Y).
\]

\[
\text{(c)}\quad
E(XY)
\leq
\sqrt{E(X^2)E(Y^2)}.
\]

\end{finalanswerbox}

\end{solutionbox}


\newpage

% ============================================================
% PROBLEM 4
% ============================================================

\begin{problemheader}
\textbf{Problem 4.} \hfill \textbf{7 points}
\end{problemheader}

Suppose
\[
X_1,\ldots,X_n
\]
are independent Bernoulli random variables with unknown success probability
\(p\), and define
\[
\widehat p
=
\frac1n\sum_{i=1}^n X_i.
\]

Let
\[
\varepsilon>0
\qquad\text{and}\qquad
0<\alpha<1.
\]

Using Chebyshev's inequality, derive a condition on the sample size \(n\) that
guarantees
\[
P\left(|\widehat p-p|<\varepsilon\right)
\geq
1-\alpha
\]
for every
\[
0<p<1.
\]

\begin{solutionbox}

Since
\[
X_i\sim\operatorname{Bernoulli}(p),
\]
we have
\[
E(X_i)=p,
\qquad
\operatorname{Var}(X_i)=p(1-p).
\]

Therefore,
\[
E(\widehat p)
=
E\left(
\frac1n\sum_{i=1}^nX_i
\right)
=
p.
\]

Since the \(X_i\)'s are independent,
\begin{align*}
\operatorname{Var}(\widehat p)
&=
\operatorname{Var}
\left(
\frac1n\sum_{i=1}^nX_i
\right)\\
&=
\frac1{n^2}
\sum_{i=1}^n
\operatorname{Var}(X_i)\\
&=
\frac{p(1-p)}{n}.
\end{align*}

Chebyshev's inequality gives
\[
P\left(
|\widehat p-p|\geq\varepsilon
\right)
\leq
\frac{\operatorname{Var}(\widehat p)}
{\varepsilon^2}
=
\frac{p(1-p)}{n\varepsilon^2}.
\]

We want a bound that holds for every \(0<p<1\).

Since
\[
p(1-p)\leq\frac14,
\]
we obtain
\[
P\left(
|\widehat p-p|\geq\varepsilon
\right)
\leq
\frac{1}{4n\varepsilon^2}.
\]

Therefore,
\[
P\left(
|\widehat p-p|<\varepsilon
\right)
\geq
1-\frac{1}{4n\varepsilon^2}.
\]

To guarantee that this is at least \(1-\alpha\), it is sufficient that
\[
\frac{1}{4n\varepsilon^2}
\leq
\alpha.
\]

Solving for \(n\),
\[
n
\geq
\frac{1}{4\alpha\varepsilon^2}.
\]

Since \(n\) must be an integer, one may take
\[
n
\geq
\left\lceil
\frac{1}{4\alpha\varepsilon^2}
\right\rceil.
\]

\begin{finalanswerbox}

A sufficient condition is
\[
n
\geq
\frac{1}{4\alpha\varepsilon^2}.
\]

Equivalently, for integer sample size,
\[
n
\geq
\left\lceil
\frac{1}{4\alpha\varepsilon^2}
\right\rceil.
\]

\end{finalanswerbox}

\end{solutionbox}


\newpage

% ============================================================
% PROBLEM 5
% ============================================================

\begin{problemheader}
\textbf{Problem 5.} \hfill \textbf{8 points}
\end{problemheader}

Let \(X\) and \(Y\) be independent exponential random variables with common
parameter \(\lambda>0\):
\[
X,Y\overset{\mathrm{i.i.d.}}{\sim}\operatorname{Exp}(\lambda),
\]
with probability density function
\[
f(w)=
\begin{cases}
\lambda e^{-\lambda w}, & w>0,\\
0, & \text{otherwise}.
\end{cases}
\]

Define
\[
T=X+Y.
\]

Using the convolution formula
\[
f_T(t)
=
\int_{-\infty}^{\infty}
f_X(x)f_Y(t-x)\,dx,
\]
find the probability density function of \(T\), including its support.

\begin{solutionbox}

Since
\[
T=X+Y
\]
and \(X,Y>0\), we must have
\[
T>0.
\]

Thus, for \(t\leq0\),
\[
f_T(t)=0.
\]

Now suppose \(t>0\).

For the integrand
\[
f_X(x)f_Y(t-x)
\]
to be nonzero, we need
\[
x>0
\]
and
\[
t-x>0.
\]

Therefore,
\[
0<x<t.
\]

Using convolution,
\begin{align*}
f_T(t)
&=
\int_0^t
f_X(x)f_Y(t-x)\,dx\\
&=
\int_0^t
\lambda e^{-\lambda x}
\lambda e^{-\lambda(t-x)}
\,dx\\
&=
\lambda^2e^{-\lambda t}
\int_0^t dx\\
&=
\lambda^2t e^{-\lambda t}.
\end{align*}

Hence
\[
f_T(t)
=
\begin{cases}
\lambda^2t e^{-\lambda t},
& t>0,\\
0,
& t\leq0.
\end{cases}
\]

Thus \(T\) does \emph{not} have an exponential distribution. It has a
Gamma distribution with shape parameter \(2\) and rate parameter
\(\lambda\).

\begin{finalanswerbox}

\[
f_T(t)
=
\begin{cases}
\lambda^2t e^{-\lambda t},
& t>0,\\
0,
& t\leq0.
\end{cases}
\]

Equivalently,
\[
T\sim\operatorname{Gamma}(2,\lambda)
\]
under the shape-rate parameterization.

\end{finalanswerbox}

\end{solutionbox}


\newpage

% ============================================================
% PROBLEM 6
% ============================================================

\begin{problemheader}
\textbf{Problem 6.} \hfill \textbf{8 points}
\end{problemheader}

Let \(X,Y,Z\) be jointly continuous with joint probability density function
\[
f_{X,Y,Z}(x,y,z)
=
\begin{cases}
x^2e^{-x(1+y+z)}, & x,y,z>0,\\
0, & \text{otherwise}.
\end{cases}
\]

The marginal probability density functions are
\[
f_X(x)
=
\begin{cases}
e^{-x}, & x>0,\\
0, & \text{otherwise},
\end{cases}
\]
and
\[
f_Y(y)
=
\begin{cases}
\dfrac{1}{(1+y)^2}, & y>0,\\
0, & \text{otherwise},
\end{cases}
\qquad
f_Z(z)
=
\begin{cases}
\dfrac{1}{(1+z)^2}, & z>0,\\
0, & \text{otherwise}.
\end{cases}
\]

\begin{enumerate}[label=(\alph*),leftmargin=*,itemsep=0.8em]

    \item Are \(X,Y,Z\) pairwise independent? Justify your answer.
    \hfill \textbf{(5 points)}

    \item Are \(X,Y,Z\) mutually independent? Justify your answer.
    \hfill \textbf{(3 points)}

\end{enumerate}

\begin{solutionbox}

\textbf{(a)} To determine pairwise independence, we examine the joint
density of each pair.

First, consider \(X\) and \(Y\). Integrating out \(Z\),
\begin{align*}
f_{X,Y}(x,y)
&=
\int_0^\infty
x^2e^{-x(1+y+z)}\,dz\\
&=
x^2e^{-x(1+y)}
\int_0^\infty e^{-xz}\,dz\\
&=
x^2e^{-x(1+y)}
\frac1x\\
&=
xe^{-x(1+y)},
\qquad x,y>0.
\end{align*}

On the other hand,
\[
f_X(x)f_Y(y)
=
e^{-x}\frac{1}{(1+y)^2}.
\]

These are not equal in general:
\[
xe^{-x(1+y)}
\neq
\frac{e^{-x}}{(1+y)^2}.
\]

Therefore, \(X\) and \(Y\) are not independent.

Consequently, \(X,Y,Z\) cannot be pairwise independent.

For completeness, the same phenomenon occurs for the other pairs.

By symmetry,
\[
f_{X,Z}(x,z)
=
xe^{-x(1+z)},
\]
which is not equal to
\[
f_X(x)f_Z(z)
=
\frac{e^{-x}}{(1+z)^2}.
\]

For \(Y\) and \(Z\),
\begin{align*}
f_{Y,Z}(y,z)
&=
\int_0^\infty
x^2e^{-x(1+y+z)}\,dx.
\end{align*}

Using
\[
\int_0^\infty x^2e^{-ax}\,dx
=
\frac{2}{a^3},
\qquad a>0,
\]
we obtain
\[
f_{Y,Z}(y,z)
=
\frac{2}{(1+y+z)^3}.
\]

However,
\[
f_Y(y)f_Z(z)
=
\frac{1}{(1+y)^2(1+z)^2},
\]
which is again different in general.

Hence none of the three pairs is independent.


\textbf{(b)} Mutual independence would require
\[
f_{X,Y,Z}(x,y,z)
=
f_X(x)f_Y(y)f_Z(z)
\]
for all points in the support.

The product of the marginals is
\[
f_X(x)f_Y(y)f_Z(z)
=
\frac{e^{-x}}
{(1+y)^2(1+z)^2}.
\]

But the joint density is
\[
f_{X,Y,Z}(x,y,z)
=
x^2e^{-x(1+y+z)}.
\]

These expressions are not equal in general.

Therefore, \(X,Y,Z\) are not mutually independent.

Also, since mutual independence implies pairwise independence and part (a)
shows that they are not pairwise independent, mutual independence is
impossible.

\begin{finalanswerbox}

\[
\text{(a)}\quad
X,Y,Z\text{ are not pairwise independent.}
\]

In fact, none of the three pairs is independent.

\[
\text{(b)}\quad
X,Y,Z\text{ are not mutually independent.}
\]

\end{finalanswerbox}

\end{solutionbox}


\newpage

% ============================================================
% PROBLEM 7
% ============================================================

\begin{problemheader}
\textbf{Problem 7.} \hfill \textbf{6 points}
\end{problemheader}

For each part below, select the \textbf{correct statement} of the indicated
theorem. Only one choice is correct in each part.

\begin{enumerate}[label=(\alph*),leftmargin=*,itemsep=1.2em]

% ------------------------------------------------------------
\item \textbf{Weak Law of Large Numbers (WLLN).}

Let \(X_1,X_2,\ldots\) be i.i.d.\ random variables with finite mean
\[
E(X_i)=\mu
\]
and finite variance
\[
\operatorname{Var}(X_i)=\sigma^2.
\]
Let
\[
\overline X_n=\frac1n\sum_{i=1}^n X_i.
\]

Which of the following is the conclusion of the Weak Law of Large Numbers?
\hfill \textbf{(2 points)}

\begin{enumerate}[label=(\Alph*),leftmargin=2em,itemsep=0.5em]

    \item For every \(\varepsilon>0\),
    \[
    P\left(
    |\overline X_n-\mu|>\varepsilon
    \right)=0
    \qquad
    \text{for every }n.
    \]

    \item For every \(\varepsilon>0\),
    \[
    \lim_{n\to\infty}
    P\left(
    |\overline X_n-\mu|>\varepsilon
    \right)
    =0.
    \]

    \item
    \[
    P\left(
    \lim_{n\to\infty}\overline X_n=\mu
    \right)=1.
    \]

    \item For every \(z\in\mathbb R\),
    \[
    \lim_{n\to\infty}
    P\left(
    \frac{\overline X_n-\mu}{\sigma/\sqrt n}
    \leq z
    \right)
    =
    \Phi(z).
    \]

\end{enumerate}


% ------------------------------------------------------------
\item \textbf{Strong Law of Large Numbers (SLLN).}

Let \(X_1,X_2,\ldots\) be i.i.d.\ random variables satisfying the same
assumptions as in the Weak Law of Large Numbers, with
\[
E(X_i)=\mu
\]
and finite variance
\[
\operatorname{Var}(X_i)=\sigma^2.
\]

Which of the following is the conclusion of the Strong Law of Large Numbers?
\hfill \textbf{(2 points)}

\begin{enumerate}[label=(\Alph*),leftmargin=2em,itemsep=0.5em]

    \item
    \[
    P\left(
    \lim_{n\to\infty}\overline X_n=\mu
    \right)=1.
    \]

    \item For every \(\varepsilon>0\),
    \[
    \lim_{n\to\infty}
    P\left(
    |\overline X_n-\mu|>\varepsilon
    \right)
    =0.
    \]

    \item
    \[
    \lim_{n\to\infty}
    P(\overline X_n=\mu)
    =1.
    \]

    \item
    \[
    \frac{\overline X_n-\mu}{\sigma/\sqrt n}
    \xrightarrow[]{d}
    N(0,1).
    \]

\end{enumerate}


% ------------------------------------------------------------
\item \textbf{Central Limit Theorem (CLT).}

Let \(X_1,X_2,\ldots\) be i.i.d.\ random variables with
\[
E(X_i)=\mu,
\qquad
\operatorname{Var}(X_i)=\sigma^2<\infty.
\]

Which of the following is the conclusion of the Central Limit Theorem?
\hfill \textbf{(2 points)}

\begin{enumerate}[label=(\Alph*),leftmargin=2em,itemsep=0.5em]

    \item
    \[
    \overline X_n
    \xrightarrow[]{P}
    N(\mu,\sigma^2).
    \]

    \item
    \[
    P\left(
    \lim_{n\to\infty}\overline X_n=\mu
    \right)=1.
    \]

    \item For every \(z\in\mathbb R\),
    \[
    \lim_{n\to\infty}
    P\left(
    \frac{\overline X_n-\mu}{\sigma/\sqrt n}
    \leq z
    \right)
    =
    \Phi(z).
    \]

    \item For every \(\varepsilon>0\),
    \[
    \lim_{n\to\infty}
    P\left(
    |\overline X_n-\mu|>\varepsilon
    \right)
    =0.
    \]

\end{enumerate}

\end{enumerate}

\begin{solutionbox}

\textbf{(a) Weak Law of Large Numbers.}

The WLLN states that the sample mean converges to the population mean
\(\mu\) \emph{in probability}. That is, for every \(\varepsilon>0\),
\[
\lim_{n\to\infty}
P\left(
|\overline X_n-\mu|>\varepsilon
\right)
=
0.
\]

Therefore, the correct answer is

\[
\textbf{B}.
\]


\textbf{(b) Strong Law of Large Numbers.}

The SLLN states that
\[
\overline X_n\longrightarrow\mu
\]
almost surely. Equivalently,
\[
P\left(
\lim_{n\to\infty}\overline X_n=\mu
\right)
=
1.
\]

Therefore, the correct answer is

\[
\textbf{A}.
\]


\textbf{(c) Central Limit Theorem.}

The CLT states that the standardized sample mean converges in distribution
to a standard normal random variable:
\[
\frac{\overline X_n-\mu}{\sigma/\sqrt n}
\xrightarrow[]{d}
N(0,1).
\]

Equivalently, for every \(z\in\mathbb R\),
\[
\lim_{n\to\infty}
P\left(
\frac{\overline X_n-\mu}{\sigma/\sqrt n}
\leq z
\right)
=
\Phi(z).
\]

Therefore, the correct answer is

\[
\textbf{C}.
\]

\begin{finalanswerbox}

\[
\boxed{\text{(a) B}}
\qquad
\boxed{\text{(b) A}}
\qquad
\boxed{\text{(c) C}}
\]

\end{finalanswerbox}

\end{solutionbox}


\end{document}